\section{A saved index needs a source version}
\label{sec:index-compatibility}

The same table and SQL primary key can survive an update that changes a
sequence. An index must therefore identify more than its owner by name. The
versioned index metadata binds the payload to a source epoch and to the rules
that produced its minimizers. Invalidation changes the old epoch's catalog
state. A replacement source receives a new epoch.

The important checks have simple reasons:
\begin{itemize}
  \item \textbf{Source identity.} The relation, column, source locator,
    epoch, and generation must name the intended source version.
  \item \textbf{Seed meaning.} The alphabet, $k,w$, hash scheme, and scheme
    version must agree with the query. Equal parameter names alone do not
    make different hash schemes compatible.
  \item \textbf{Position meaning.} Width and coordinate encoding must
    represent the target without truncation, including its strand bits.
  \item \textbf{Array meaning.} The reader must recognize the format and
    every addressed range. A device reader also needs a compatible array
    interface.
\end{itemize}

Lengths, offsets, and checksums help detect damaged or misaddressed data.
Checks happen before the reader forms a span or allocates a destination from
an untrusted stored count. A checksum is a corruption check; it does not
replace source-version or algorithm checks. Ordinary serialized minimizer
formats record layout revisions and array sizes, but do not store a source
epoch or hash-scheme identifier. Generation readers also check the stored
source version and minimizer hash scheme.

\subsection{Publish complete generations}
\label{sec:index-generations}
\label{sec:index-physical-persistence}

A small \emph{descriptor} is the record that identifies a generation and its
payload. Building a new generation does not require changing
the old payload in place:
\begin{enumerate}
  \item Allocate a new generation and write its payload.
  \item Check the complete layout and its source and algorithm identities.
  \item Replace the catalog descriptor transactionally, after the payload
    is ready.
\end{enumerate}
The descriptor update uses PostgreSQL's tuple visibility. An older snapshot
can continue to see its older descriptor; a later eligible snapshot can see
the replacement. Failure before publication must not make a partial build
the current index.

\begin{figure}[htbp]
\centering
\begin{tikzpicture}[x=1cm,y=1cm,font=\footnotesize]
  \draw[index arrow,black!50] (1,5.6)--(12.2,5.6) node[right] {time};
  \node[anchor=east] at (.8,4.55) {writer};
  \node[index orange,minimum width=4.2cm] at (4.1,4.55) {write and check generation $g+1$};
  \node[index green,minimum width=2.8cm] at (9.05,4.55) {publish; commit};
  \draw[index arrow] (6.3,4.55)--(7.6,4.55);
  \draw[dashed,black!60] (9.05,.85)--(9.05,4.2);
  \draw[dashed,black!60] (9.05,4.9)--(9.05,5.2);
  \node[anchor=east] at (.8,3.05) {reader A};
  \draw[line width=8pt,indexBlue!25] (1.8,3.05)--(10.9,3.05);
  \node at (5.45,3.05) {older snapshot: descriptor for $g$};
  \node[anchor=east] at (.8,1.6) {reader B};
  \draw[line width=8pt,indexGreen!25] (9.45,1.6)--(12,1.6);
  \node[anchor=east] at (8.85,1.6) {new snapshot};
  \node at (10.75,1.6) {$g+1$};
  \node[anchor=west,align=left] at (1.6,.3)
    {Publication changes which descriptor is visible.\\It does not overwrite the payload still used by reader A.};
\end{tikzpicture}
\caption{Versioned publication with two snapshots. Reader A resolves the
old descriptor and retains its read resources. Reader B starts after commit
and can resolve the new one. Physical reclamation is a separate operation.}
\label{fig:index-generation-lifecycle}
\end{figure}

A snapshot that sees a parent sequence's deletion must also reject its old
source epoch. These changes belong to the same transaction. Updating a
sequence similarly invalidates the old source for newer snapshots and
appends a replacement. Neither
operation means that every fragment and index byte is immediately erased.
An old reader may still need the former version, and the stored objects need
their own cleanup policy. The current MPS build path does not automatically
delete replaced payloads; reclaiming them requires separate cleanup.

\subsection{From stored index rows to a CPU reader file}
\label{sec:index-mpf5-sidecar}
\label{sec:index-formats}

The optional grouped MPF path keeps two stored copies of its index data.
It first writes index rows in PostgreSQL, then creates a file for CPU reads.
To understand why, distinguish three objects:
\begin{itemize}
  \item \textbf{The sequence.} The original packed symbols remain in the
    sequence-storage relations from the previous chapter.
  \item \textbf{The index relation.} Separate heap rows hold the computed
    minimizer controls and positions, split into byte chunks, with generation
    metadata. This is already an index, not another copy of the sequence.
  \item \textbf{The reader file.} This file contains the index payload bytes
    in consecutive ranges, their directory and metadata, and an additional
    direct hash lookup. We call this fixed layout a \emph{reader image}.
    Once published, the image is read-only.
\end{itemize}

The conversion copies the encoded index payload out of the rows. It does
not copy PostgreSQL pages or recompute the minimizers. Heap pages contain
page headers, row headers and other fields around each stored value
\autocite{postgresql17_storage_page_layout}. Mapping the heap file would
still leave those fields between payload chunks. It would not expose one
contiguous index array. The reader image removes that obstacle once during
index construction.

\paragraph{Construction happens once per generation.}
The current relation-backed builder takes these steps:
\begin{enumerate}
  \item \textbf{Build the index rows.} Compute the minimizers and save a new
    generation's payload, directory and metadata in the index relation.
  \item \textbf{Write the reader file.} Reread and check that generation's
    rows in index order. Append their payload bytes to a private construction
    file. Add the directory, metadata and direct hash lookup, then finish
    and check the file header.
  \item \textbf{Finish the file before publishing the index.} Make the file
    read-only, synchronize it, give it a new final name without replacing an
    existing file, and remove the temporary name. Then update the database
    descriptor and active-generation record in the transaction. Other
    transactions can select the new generation after commit.
\end{enumerate}
The index relation remains after this export. The extra copy costs disk space
and a complete payload reread during construction. The direct hash lookup
also adds construction work and bytes. In the current writer this section is
always built; the lookup setting chooses how a later query reads the image.
Figure~\ref{fig:index-control-data-plane} shows generation $g$ and its
completed file $F_g$.

\begin{figure}[htbp]
\centering
\begin{tikzpicture}[x=1cm,y=1cm,font=\footnotesize]
  \node[anchor=west,font=\bfseries] at (0,7.0)
    {(a) During construction of generation $g$};
  \node[anchor=west] at (.3,6.35) {Index relation: three payload chunks};
  \node[anchor=west] at (7.0,6.35) {Completed reader file $F_g$};
  \draw[rounded corners,draw=black!55] (6.9,2.85) rectangle (12.2,6.05);
  \foreach \y/\n/\col in {5.65/0/indexBlue,4.85/1/indexOrange,4.05/2/indexGreen} {
    \draw[fill=black!5,draw=black!55] (.3,\y-.3) rectangle (1.9,\y+.3);
    \node at (1.1,\y) {row fields};
    \draw[fill=\col!15,draw=\col] (1.9,\y-.3) rectangle (4.65,\y+.3);
    \node at (3.275,\y) {chunk \n};
  }
  \draw[fill=black!5,draw=black!55] (7.0,5.35) rectangle (12.1,5.95);
  \node at (9.55,5.65) {header: identity and byte ranges};
  \foreach \x/\n/\col in {7.0/0/indexBlue,8.7/1/indexOrange,10.4/2/indexGreen} {
    \draw[fill=\col!15,draw=\col] (\x,4.55) rectangle (\x+1.7,5.15);
    \node at (\x+.85,4.85) {chunk \n};
  }
  \draw[fill=black!5,draw=black!55] (7.0,3.75) rectangle (12.1,4.35);
  \node at (9.55,4.05) {directory and generation metadata};
  \draw[fill=black!5,draw=black!55] (7.0,2.95) rectangle (12.1,3.55);
  \node at (9.55,3.25) {direct hash lookup};
  \draw[index arrow] (4.85,4.85)--(6.8,4.85);
  \node[align=center] at (5.82,4.05) {copy payload\\bytes only};
  \node[anchor=west,font=\bfseries] at (0,2.3)
    {(b) A later query: open the existing file};
  \node[index blue,minimum width=4.5cm,minimum height=.85cm] (desc)
    at (2.55,1.35) {Visible database descriptor\\source version and $g$};
  \node[anchor=west] at (7.0,1.85) {File pages mapped for this query};
  \foreach \i in {0,...,5}
    \draw[fill=black!4,draw=black!55]
      (7.0+.85*\i,1.05) rectangle (7.85+.85*\i,1.65);
  \foreach \i in {0,4}
    \draw[fill=indexGreen!20,draw=indexGreen]
      (7.0+.85*\i,1.05) rectangle (7.85+.85*\i,1.65);
  \draw[index arrow] (desc.east)--(6.8,1.35);
  \node[align=center] at (5.85,.62) {select and\\check $F_g$};
  \draw[index arrow] (12.2,4.45)--(12.65,4.45)--(12.65,1.35)--(12.3,1.35);
  \node[anchor=west,align=left] at (.3,-.1)
    {The query does not repeat the export.};
  \node[anchor=west,align=left] at (7.0,-.1)
    {Green: pages used by this lookup.\\The full file need not be resident.};
\end{tikzpicture}
\caption{The same index payload has two storage layouts. Colors trace
three illustrative chunks from heap rows into consecutive file ranges.
Row fields are left behind. A later query checks the visible generation and
maps its completed image. The lower strip shows file pages, not index-array
boundaries; the operating system may also read pages ahead of a request.}
\label{fig:index-control-data-plane}
\end{figure}

\paragraph{Queries open the file rather than exporting it again.}
The CPU reader first validates the source and reads the generation descriptor
under the query snapshot. It opens the matching file read-only, checks its
identity and size, and maps it with \texttt{mmap}. Mapping gives the process
an address range through which it can read the file. After header and range
checks, the reader uses views of the mapped index arrays. It still allocates
metadata and query work buffers, but does not copy the complete index
payload into a new array. One prepared reader can serve several query
sequences in a batch. A later prepared reader can reopen the same file.

Algorithm~\ref{alg:index-image-lifecycle} separates the two operations.
Here $s$ is the complete, unchanged stored target sequence, $R_g$ denotes
its index rows for generation $g$, and $\sigma$ is the query snapshot.
The named operations are pseudocode, not an API.
File and database errors stop the operation; registered cleanup releases
partly opened resources. The build runs within the caller's transaction.

\begin{algorithm}[htbp]
\caption{Build an MPF reader image once; open it for later queries.}
\label{alg:index-image-lifecycle}
\small
\begin{algorithmic}[1]
  \Procedure{BuildMappedMPF}{$s$}
    \State $g\gets$ a new generation number
    \State $R_g\gets\Call{BuildAndStoreMPF}{s,g}$
    \State $F\gets\Call{CreatePrivateFile}{}$
    \ForAll{payload chunks $c$ of $R_g$, in stored index order}
      \State \Call{CheckChunk}{$c$}
      \State \Call{AppendPayload}{$F,c$}
    \EndFor
    \State \Call{AddMetadataAndDirectLookup}{$F,R_g$}
    \State $F_g\gets\Call{CheckAndPublishReadOnlyFile}{F,g}$
    \State \Call{UpdateDescriptorAndActiveGeneration}{$s,g,F_g$}
    \State \Return $g$ \Comment{the caller commits the transaction}
  \EndProcedure
  \Statex
  \Procedure{OpenMappedMPF}{$s,\sigma$}
    \State $D\gets\Call{ReadAndCheckVisibleDescriptor}{s,\sigma}$
    \State $F_g\gets\Call{OpenAndCheckFileReadOnly}{D}$ \Comment{check identity and size}
    \State $M\gets\Call{MapFileReadOnly}{F_g}$
    \State \Call{CheckHeaderAndRanges}{$M,D$}
    \State \Return a reader using $M$ while the read session owns the mapping and file
  \EndProcedure
\end{algorithmic}
\end{algorithm}

\paragraph{This is an optional SQL path.}
The public function \texttt{bioseq.update\_index} reaches the file writer
when grouped MPF construction is enabled. It is not limited to a test fixture.
The mapped CPU reader runs in \texttt{BioseqBatchJoin}.
Set the per-call \texttt{index\_read\_requirement} option to 2,
meaning \emph{require a compatible stored index}. Ordinary column scans
do not support that requirement, and the default permissive search does not
select this mapped route. Enabling grouped construction alone does not
select it for later queries.

Listing~\ref{lst:index-mapped-user-path} shows the public calls. It assumes
committed DNA4 sequences with 32-bit packed units in a new permanent
\texttt{refs(id,seq)} table using the Bioseq table access method, and query
sequences with the same alphabet and packing in \texttt{queries(id,seq)}.
Load \texttt{refs} with
\texttt{bioseq.index.insert\_maintenance.enabled = off}, so it has no older
ordinary MPF relation. The example
uses the registered $k=15,w=10$ profile. The same minimizer algorithm is kept
for construction and search.

\noindent\begin{minipage}{\linewidth}
\begin{lstlisting}[style=bioseq-index-sql,
  caption={Explicitly build and use the mapped MPF path.},
  label={lst:index-mapped-user-path}]
BEGIN;
SET max_parallel_workers_per_gather = 0;
SET bioseq.aligner.vulkan.enabled = off;
SET bioseq.aligner.minimizer.algorithm = 'two_stack_nthash';
SET bioseq.aligner.minimizer.kmer_size = 15;
SET bioseq.aligner.minimizer.window_size = 10;
SET bioseq.aligner.chaining.algorithm = 'simd_minimap2';
SET bioseq.aligner.base_alignment.algorithm = 'block_aligner';
SET bioseq.index.async_persistence.enabled = off;
SET bioseq.index.recompute_async_persistence.enabled = off;
SET bioseq.aligner.grouped_soa_generation.enabled = on;

SELECT bioseq.update_index('refs'::regclass,
                          'minimizers_per_fragment');
COMMIT;

SET bioseq.index.sequence_read.enabled = on;
SELECT bioseq.set_active_read_index('refs'::regclass,
                                   'minimizers_per_fragment');
SELECT r.id, q.id
FROM refs AS r JOIN queries AS q
  ON mmsearch(r.seq, q.seq,
              aligner_options(index_read_requirement => 2));
\end{lstlisting}
\end{minipage}

The default grouped MPF control reader and the optional direct hash reader
both use the image. MPS instead reads its relation-backed streams. The
grouped MPF GPU path reads the complete relation payload into device arrays.
These are separate read paths.

\paragraph{RAM is an alternative, not the opposite of a mapped file.}
An index kept in RAM can be fast while it remains available. BioSeqDB's
ordinary index caches already retain prepared CPU objects. However, a new
backend may have to read and reconstruct its own copy, and an index may be
larger than the memory available for that copy.

A mapped file still uses RAM: the operating system keeps file pages in its
page cache. Read-only mappings in different processes can use the same
cached file pages. Pages can be loaded as needed and reclaimed under memory
pressure without making each backend retain the whole index in a private
array. This is the performance reason for the extra representation: avoid
repeated tuple extraction and full-payload copies at query preparation.
It is not a guarantee of faster search. A cold query can wait for page reads,
and a broad lookup can touch much of the file. Building the image also costs
an extra write, checks and synchronization. A warm in-memory index remains
a useful comparison.

\paragraph{Releasing a reader and deleting its file are different steps.}
The file has the following lifecycle:
\begin{itemize}
  \item \textbf{During construction:} a private temporary name identifies
    the unfinished file. Publication installs a distinct final name.
  \item \textbf{On transaction abort:} registered file actions attempt to remove the
    transaction's unfinished or newly published files after checking their
    identity. On commit, the completed file is kept.
  \item \textbf{When a reader finishes:} its PostgreSQL resource owner
    unmaps the address range and closes the file descriptor. The file remains
    available for another reader. Error cleanup releases these resources too.
  \item \textbf{After replacement or source deletion:} visibility and source
    checks decide whether a reader may use the old generation. The current
    query and rebuild paths do not automatically reclaim its retired image.
\end{itemize}

A missing or damaged image causes the required mapped read to fail. The
query does not recreate it from the relation. An explicit rebuild produces a
new generation and image. The file is derived data outside ordinary relation
WAL. Synchronizing it does not establish a recovery protocol or ensure that
a backup contains a matching file and database descriptor.
Its relation-backed generation also uses unlogged payload
relations. Crash recovery, backup and removal of orphaned files therefore
need explicit support rather than an assumption that PostgreSQL manages the
image like a heap relation.
